%%%%%
%
%% Dans ce fichier, les �ventuels accents sont cod�s
%% en ISO-8859-15 (aka latin9)
%% La ligne suivante donne trois exemples 
%% e accent aigu : �, e accent grave : �, e-dans-l'o : �
\section{Polyn�me de 5$^{eme}$ degr� appliqu� � un param�tre \texttt{float}}

\itemstyle{Op\'erateur~:} \operfmt{Poly\_5\_flt}
\index{Poly\_5\_flt}

%%%%%OBJET%%%%
\itemstyle{Objet~:}
Application d'un polyn�me de 5$^{eme}$ degr�
� un param�tre.

%%%%%ARGUMENTS%%%%
\bigskip
%% Dans ce fichier, les �ventuels accents sont cod�s
%% en ISO-8859-15 (aka latin9)
%% La ligne suivante donne trois exemples 
%% e accent aigu : �, e accent grave : �, e-dans-l'o : �
%%%%%ENTREES%%%%
\begin{tabarg}
\multicolumn{3}{l}{\titretab {Entr\'ees}}\\
\hline
$X$ & Vecteur,  \texttt{float} & Op�rande (unit� quelconque)\\
$C_0$ & Coeff., \texttt{float} & Coefficient du terme de degr� 0\\
$C_1$ & Coeff., \texttt{float} & Coefficient du terme de degr� 1\\
$C_2$ & Coeff., \texttt{float} & Coefficient du terme de degr� 2\\
$C_3$ & Coeff., \texttt{float} & Coefficient du terme de degr� 3\\
$C_4$ & Coeff., \texttt{float} & Coefficient du terme de degr� 4\\
$C_5$ & Coeff., \texttt{float} & Coefficient du terme de degr� 5\\
$Unit\_res$ & Coeff., \texttt{String} & Unit� du r�sultat\\
\end{tabarg}

\medskip
%%%%%SORTIES%%%%
\begin{tabarg}
\multicolumn{3}{l}{\titretab {Sorties}}\\
\hline
$Z$ & Vecteur, \texttt{float} & R�sultat (unit� indiqu�e par $Unit\_res$)\\
\end{tabarg}


%%%%%FORMULE%%%%
\itemstyle{Formule}

%% Formule de l'operateur : Poly_5_flt
%% Dans ce fichier, les �ventuels accents sont cod�s
%% en ISO-8859-15 (aka latin9)
%% La ligne suivante donne trois exemples 
%% e accent aigu : �, e accent grave : �, e-dans-l'o : �
Le polyn�me $ Z = C_0 + C_1 . X + C_2 .X^2 + C_3 . X^3 + C_4 .X^4 + C_5 . X^5$
est calcul� par la m�thode dite "de Horner" :
\[
Z = C_0 + X . (C_1 + X . ( C_2 + X . ( C_3 + X . ( C_4 + X . C_5))))
\]



%%%%%AUTEUR%%%%
\itemstyle{Auteur~:}
CNRM/GMEI/TRAMM

%%%%%REFERENCE%%%%
\itemstyle{R\'ef\'erences~:}
La m�thode "de Horner" est r�f�renc�e dans TAOCP \cite{knuth-sa-97}, section
4.6.4, pp 485-487, qui explique qu'elle �tait d�j� connue de Newton
et du math�maticien chinois Ch'in Chiu Shao (13eme si�cle).
